\documentclass{template}

\usepackage[square,numbers]{natbib}

\usepackage{hyperref,xcolor}
\hypersetup{
    pdftitle={CV - Alexander Weichart},
    pdfauthor={Alexander Weichart},
    pdfsubject={Full-Stack Software Engineer - AI-Native Development - Resume},
    pdfkeywords={Alexander Weichart, CV, Resume, Software Engineer, Full Stack, Java, TypeScript, Spring Boot, React, PostgreSQL, Kubernetes, GitLab, Flux CD, Linux, DevOps, Agile, AI, LLM, Claude Code, Codex, Copilot, MCP, Agentic Engineering},
    pdfcopyright={Copyright (C) 2026 Alexander Weichart},
    pdfmetalang={en},
    linkcolor=blue,
    citecolor=blue
}

\setname{Alexander Weichart}{}
\setposition{Full-Stack Software Engineer | AI-Native Development}
\setaddress{Osaka, Japan}
\setmobile{\getPhone}
\setmail{alexanderweichart@fastmail.com}
\setcitizenship{German}
\setlinkedinaccount{https://www.linkedin.com/in/alexander-weichart-a357241a4/} 
\setprojects{https://alexanderweichart.de/4_Projects/Projects} 
\sethomepage{https://www.alexanderweichart.de/}
\setgithubaccount{https://github.com/AlexW00} 
\setprofilepic{profile.jpeg}
\setthemecolor{Black}


\begin{document}

\headerview
\vspace{3ex}

%
\section{Skills}
    \newcommand{\skillone}{\createskill{AI Tools}{Claude Code \cpshalf Codex \cpshalf GitHub Copilot}}
    \newcommand{\skilltwo}{\createskill{Agentic Engineering}{Skills, sub-agents, MCP \cpshalf Context engineering, model selection}}
    \newcommand{\skillthree}{\createskill{Core Languages}{Java \cpshalf TypeScript \cpshalf Swift}}
    \newcommand{\skillfour}{\createskill{Core Technologies}{Spring Boot, React, PostgreSQL \cpshalf Kubernetes, GitLab, Flux \cpshalf Linux}}
    \newcommand{\skillfive}{\createskill{Languages}{German (native) \cpshalf English (fluent) \cpshalf Japanese (conversational) \cpshalf French (basic)}}
    %
    \createskills{\skillone, \skilltwo, \skillthree, \skillfour, \skillfive}
    \vspace{-3mm} % reduces whitespace!
    %
\bigskip

    \section{Experience}
    %
    \datedexperience{Birchill (Osaka, Japan)}{Apr 2026 - Present}
    \explanation{Full-Stack Software Engineer (Working Holiday) - TypeScript, AWS}{Osaka, Japan}
    \explanationdetail{
    \smallskip
    \coloredbullet\ %
    Lead end-to-end development of a full-stack Japanese learning platform, owning technical direction and delivery across the frontend, backend, and AWS infrastructure.

    \smallskip
    \coloredbullet\ %
    Build TypeScript services and AWS workflows with Lambda, DynamoDB, Step Functions, and CDK for \mbox{dictionary data} and learning features.
    }
    %
    \datedexperience{DEBTVISION GmbH (Baden-Württemberg State Bank)}{Oct 2023 - Mar 2026}
    \explanation{Full-Stack Software Engineer - Java, Spring Boot, TypeScript, Kubernetes, GitLab}{Stuttgart (remote), Germany}
    \explanationdetail{
    \smallskip
    \coloredbullet\ %
    Owned features end-to-end on a B2B financial platform, spanning API design, database modeling, CI/CD pipelines, and production deployments on Kubernetes.

    \smallskip
    \coloredbullet\ %
    Represented engineering in a cross-functional AI initiative; introduced GitHub Copilot as the team's first AI coding tool and led adoption of agentic workflows through workshops, mentoring, and workflow design, contributing to an 80\% increase in \mbox{sprint velocity} over six months.

    \smallskip
    \coloredbullet\ %
    Right-sized cloud infrastructure to double CPU and increase memory by 12\%, while cutting costs by 17\%.

    \smallskip
    \coloredbullet\ %
    Implemented automated tooling to restore the production cluster as a separate recovery cluster to a different cloud provider from scratch.
    }
    %
    \datedexperience{University of Regensburg}{Sep 2020 - Sep 2023}
    \explanation{Student Research Assistant - JavaScript, Kotlin, C, Neo4j, Figma}{Regensburg, Germany}
    \explanationdetail{
    \smallskip
    \coloredbullet\ %
    Led the development of \href{https://hci.ur.de/projects/graphit}{GraphIT}, a platform for organizing curricula as graphs, published at DELFI 2025.
    }

\bigskip

\section{\texorpdfstring{Projects\ \projectslinkicon}{Projects}}
    %
    \datedexperience{Weichart Apps\ \href{https://apps.weichart.de}{{\normalfont\mdseries\upshape\textcolor{NavyBlue}{\raisebox{0.15ex}{\faLink}}}}}{}
    \explanationdetail{
    \smallskip
    \coloredbullet\ %
    Designed and shipped \href{https://apps.weichart.de}{three carefully crafted iOS apps}, including \href{https://apps.apple.com/app/id6747954401}{Hibi}, a paper-like calendar, reaching over 15k downloads; built with agentic development workflows.
    }
    %
    \datedexperience{AI-Native Knowledge System\ \href{https://alexanderweichart.de/4_Projects/obsidian-setup-2026/My-AI-native-Obsidian-Setup}{{\normalfont\mdseries\upshape\textcolor{NavyBlue}{\raisebox{0.15ex}{\faLink}}}}}{}
    \explanationdetail{
    \smallskip
    \coloredbullet\ %
    Built a Git-backed knowledge system with scheduled Claude Code routines, reusable skills, MCP/CLI integrations, persistent state, and approval gates; the public site receives over 3k visits per month.
    }
    %
    \datedexperience{Obsidian 3D Graph\ \href{https://github.com/AlexW00/obsidian-3d-graph}{{\normalfont\mdseries\upshape\textcolor{NavyBlue}{\raisebox{0.15ex}{\faLink}}}}}{}
    \explanationdetail{
    \smallskip
    \coloredbullet\ %
    Obsidian plugin to visualize notes as a 3D WebGL graph, counting over 50k downloads and 350 stars on GitHub.
    }
    %
    \datedexperience{AI Prototypes (since 2023)\ \href{https://alexanderweichart.de/4_Projects/Projects}{{\normalfont\mdseries\upshape\textcolor{NavyBlue}{\raisebox{0.15ex}{\faLink}}}}}{}
    \explanationdetail{
    \smallskip
    \coloredbullet\ %
    Building with LLMs since 2023: \href{https://alexanderweichart.de/4_Projects/Projects\#-gpt-lang-2023-03-24}{GPT-Lang}, a natural-language-to-JavaScript compiler; \href{https://alexanderweichart.de/4_Projects/Projects\#-tandem-gpt-2023-02-04}{Tandem GPT}; and a recent \href{https://alexanderweichart.de/4_Projects/Projects\#ai-hud-for-japanese-text-2025-12-22}{AI-HUD for Japanese text recognition}.
    }

\bigskip

\section{Education}
    \datedexperience{Bachelor of Arts - Media Computer Science (GPA: 3.5 / 4.0)}{Sep 2019 - Sep 2023}
    \explanation{University of Regensburg}{Regensburg, Germany}
    \explanationdetail{
        \smallskip
        \coloredbullet\ %
        Full Stack Development, Agile Project Management, Natural Language Processing, Ethics.
     }

\bigskip

\section{Community}
\indent
\explanationdetail{
\smallskip
\coloredbullet\ %
Maintainer of multiple popular open-source projects (\href{https://github.com/AlexW00}{GitHub}), counting more than 1k stars in total.

\smallskip
\coloredbullet\ %
Supporting member at the German-Japanese Society Regensburg (\href{https://djg-regensburg.de/de/}{DJG Regensburg}).

}
\bigskip

\section{Publications}
\indent
\explanationdetail{
\smallskip
\coloredbullet\ %
Andreas S., Thomas F., Alexander W., Alexander H., and Raphael W., ``ScreenshotMatcher: Taking Smartphone Photos to Capture Screenshots.'' \emph{MuC '21: Proceedings of Mensch und Computer 2021}, Ingolstadt, Germany.

\smallskip
\coloredbullet\ %
Raphael W., Leonie S., and Alexander W., ``Demonstrating GraphIT - a Platform for Working with Dependency Graphs of Learning Topics.'' \emph{23. Fachtagung Bildungstechnologien (DELFI 2025)} (\href{https://dl.gi.de/handle/20.500.12116/47137}{GI DL}).
 }

\end{document}
